\documentclass[11pt]{article}
\usepackage[UTF8]{ctex}
\usepackage{amsmath,amssymb,amsthm}
\usepackage{etoolbox}

% 1. 自定义定理样式实现标题后强制换行
\newtheoremstyle{break}% 样式名称
  {\topsep}{\topsep}% 环境上下间距
  {\normalfont}% 正文字体
  {}% 缩进量
  {\bfseries}{.}% 头部字体和标点
  {\newline}% 头部后的间距：设置为 \newline 即可实现完美的换行效果
  {}% 头部时间/附加参数规范

% 应用新样式并声明常用的定理类环境
\theoremstyle{break}
\newtheorem{theorem}{定理}[section]
\newtheorem{lemma}[theorem]{引理}
\newtheorem{corollary}[theorem]{推论}
\newtheorem{definition}[theorem]{定义}

% 2. 利用 etoolbox 宏包动态修补 proof 环境，使其标题后自动换行
\patchcmd{\proof}{\item[\hskip\labelsep\itshape#1\@addpunct{.}]}{\item[\hskip\labelsep\itshape#1\@addpunct{.}]\hfill\par}{}{}

\begin{document}
\section{摘要}
对于固定的 $h$，令 $N(h,k)$ 为满足如下条件的最小正整数 $N$：任意由 $k$ 个整数组成的集合所达到的基数 $|hA|$，亦必可由 $\{1,\dots,N\}$ 的某一 $k$ 元子集达到。借助 Rajagopal 关于固定 $h$ 情形下 $|hA|$ 取值范围的定理，我们证明了对任意固定的 $h$，均有 $N(h,k) \leq k^{C_h}$。本文的核心新贡献在于，将 Rajagopal 大值构造中所用的稀疏几何块替换为具有多项式直径的结构。该替换操作在低次 Hilbert 函数层面实施，而由于当 $h$ 固定时，低次 Hilbert 函数是刻画 $h$ 重集和的唯一相关数据，故该替换足以完成论证。
\section{Introduction}\n对于有限集 $A \subset \mathbb{Z}$ 与正整数 $h$，记\[ hA = \left\{ \sum_{i=1}^h a_i : a_i \in A \right\} \]\n其中允许元素重复出现。\n定义\[ R(h,k) = \{ |hA| : A \subset \mathbb{Z}, |A| = k \}. \]\n令 $N(h,k)$ 为满足以下条件的最小正整数 $N$：\[ R(h,k) = \{ |hA| : A \subseteq \{1,\dots,N\}, |A| = k \}. \]\n等价地，我们亦可在区间 $[0, N-1]$ 中进行讨论，并于末尾施加平移变换。\n本文旨在证明下述多项式压缩界。
\textbf{定理 1.1.} 对于任意固定的 $h \geq 1$，存在常数 $C$，使得对所有 $k \geq 2$，有
\[
N(h,k) \leq k^{Ch}.
\]
该证明基于 Rajagopal 关于固定折叠和集大小取值范围的结果。令
\[
M = h^{k-h+1}, \quad K = \binom{h+k-1}{h,k},
\]
并定义
\[
\Delta = \prod_{\ell=0}^{\min(h,k)-3} [M_{h,k} + \ell h + 1, M_{h,k} + \ell h + (h-2-\ell)].
\]
此处及下文，$[u,v]$ 表示满足 $u \leq n \leq v$ 的整数 $n$ 的集合；若 $u > v$，则该区间为空集。
\end{document}